% Source: https://artofproblemsolving.com/community/c3h2304256_solution_to_problem_18_of_the_amc_8
% Stable identity: aops-topic-2304256-q01
Each row and each column of the $5\times5$ array below is an arithmetic sequence with five terms. The four corner entries are shown, and the centre entry is $X$.
\[\begin{array}{|c|c|c|c|c|}\hline
1&\phantom{00}&\phantom{00}&\phantom{00}&25\\\hline
\vphantom{1}&&&&\\\hline
&&X&&\\\hline
\vphantom{1}&&&&\\\hline
17&&&&81\\\hline
\end{array}\]
\begin{enumerate}
\item[(a)] What is the value of $X$?

$\textbf{(A) }21\qquad\textbf{(B) }31\qquad\textbf{(C) }36\qquad\textbf{(D) }40\qquad\textbf{(E) }42$
\item[(b)] A student computes $X=\dfrac12\left(\dfrac{1+81}{2}+\dfrac{17+25}{2}\right)$, the average of the averages of the corner entries on the two diagonals. Explain why this gives the right value, and show that $X$ is equal to neither $\dfrac{1+81}{2}$ nor $\dfrac{17+25}{2}$, so that neither diagonal is itself an arithmetic sequence.
\end{enumerate}
